% GENERATED FILE. Edit canonical lessons, then run npm run generate:books.
% canonical-source-hash: fnv1a64:9609845e965240f2
% canonical-lessons: RU-C71-sunday, RU-C71-winter, RU-C71-spring, RU-C71-summer, RU-C71-autumn

\chapter{Sunday and the Four Seasons}
\label{ch:ru-sunday-and-the-four-seasons}
\hlchaptermodality{71}

\begin{chapteropening}
\textbf{By the end of this chapter:} \emph{I can name Sunday, winter, spring, summer and autumn.}

Complete the last lesson of chapter 71: I can name Sunday, winter, spring, summer and autumn.
\end{chapteropening}

\section[\texorpdfstring{\ru{воскресенье} (voskresénye)}{воскресенье (voskresénye)}]{\ru{воскресенье} (voskresénye) — Sunday}

\label{lesson:RU-C71-sunday}



\begin{quote}
Before the new one: say the Russian for Friday, then the Russian for Saturday.
\end{quote}

\subsection*{You'll want to know: \ru{воскресенье}}
\textbf{\ru{воскресенье}} (\emph{voskresénye}) — \textquotedblleft{}Sunday\textquotedblright{}.

\begin{cousinweb}
\textbf{\ru{воскресенье}} — Church Slavonic, says Vasmer, as its \emph{\ru{вос}-} shows: it meant \textquotedblleft{}resurrection\textquotedblright{}, then \textquotedblleft{}the day of the Resurrection\textquotedblright{}, then Sunday. Earlier, \emph{\ru{неделя}} named Sunday.
\end{cousinweb}

\begin{grammarlens}[title={What you've built}]
One more word for this chapter.
\end{grammarlens}

\subsection*{Your turn}
\begin{itemize}
\raggedright
  \item \emph{Say it:} \emph{voskresénye}
  \item \emph{Say it:} \emph{voskresénye}, once more
  \item \emph{Say it:} say \emph{voskresénye}
  \item \emph{Recall:} say the Russian for Friday, then the Russian for Saturday, then say \emph{voskresénye} again
\end{itemize}

\begin{tcolorbox}[breakable,colback=teal!4,colframe=teal!35!black,title={Before you move on}]
What does \ru{воскресенье} mean? (\textquotedblleft{}Sunday\textquotedblright{}.) Say it once more.
\end{tcolorbox}
\section[\texorpdfstring{\ru{зима} (zimá)}{зима (zimá)}]{\ru{зима} (zimá) — winter}

\label{lesson:RU-C71-winter}



\begin{quote}
Before the new one: say the Russian for Saturday, then the Russian for Sunday.
\end{quote}

\subsection*{You'll want to know: \ru{зима}}
\textbf{\ru{зима}} (\emph{zimá}) — \textquotedblleft{}winter\textquotedblright{}.

\begin{cousinweb}
\textbf{\ru{зима}} $\leftarrow$ Proto-Slavic \emph{*zima} — inherited. Vasmer lists Lithuanian \emph{žiemà}, Sanskrit \emph{himá-}, Greek \emph{cheîma} and Latin \emph{hiems}, all \textquotedblleft{}winter\textquotedblright{}, from PIE \emph{*ǵ\textsuperscript{h}éyōm}, \textquotedblleft{}winter\textquotedblright{}.
\end{cousinweb}

\begin{grammarlens}[title={What you've built}]
One more word for this chapter.
\end{grammarlens}

\subsection*{Your turn}
\begin{itemize}
\raggedright
  \item \emph{Say it:} \emph{zimá}
  \item \emph{Say it:} \emph{zimá}, once more
  \item \emph{Say it:} say \emph{zimá}
  \item \emph{Recall:} say the Russian for Saturday, then the Russian for Sunday, then say \emph{zimá} again
\end{itemize}

\begin{tcolorbox}[breakable,colback=teal!4,colframe=teal!35!black,title={Before you move on}]
What does \ru{зима} mean? (\textquotedblleft{}Winter\textquotedblright{}.) Say it once more.
\end{tcolorbox}
\section[\texorpdfstring{\ru{весна} (vesná)}{весна (vesná)}]{\ru{весна} (vesná) — spring}

\label{lesson:RU-C71-spring}



\begin{quote}
Before the new one: say the Russian for Sunday, then the Russian for winter.
\end{quote}

\subsection*{You'll want to know: \ru{весна}}
\textbf{\ru{весна}} (\emph{vesná}) — \textquotedblleft{}spring\textquotedblright{}.

\begin{cousinweb}
\textbf{\ru{весна}} $\leftarrow$ Proto-Slavic \emph{*vesna} — inherited. Vasmer lists Lithuanian \emph{vasarà}, \textquotedblleft{}summer\textquotedblright{}, Greek \emph{éar}, Latin \emph{vēr} and Old Norse \emph{vár}, all \textquotedblleft{}spring\textquotedblright{}, and Sanskrit \emph{vasantá-}, \textquotedblleft{}spring\textquotedblright{}.
\end{cousinweb}

\begin{grammarlens}[title={What you've built}]
One more word for this chapter.
\end{grammarlens}

\subsection*{Your turn}
\begin{itemize}
\raggedright
  \item \emph{Say it:} \emph{vesná}
  \item \emph{Say it:} \emph{vesná}, once more
  \item \emph{Say it:} say \emph{vesná}
  \item \emph{Recall:} say the Russian for Sunday, then the Russian for winter, then say \emph{vesná} again
\end{itemize}

\begin{tcolorbox}[breakable,colback=teal!4,colframe=teal!35!black,title={Before you move on}]
What does \ru{весна} mean? (\textquotedblleft{}Spring\textquotedblright{}.) Say it once more.
\end{tcolorbox}
\section[\texorpdfstring{\ru{лето} (léto)}{лето (léto)}]{\ru{лето} (léto) — summer}

\label{lesson:RU-C71-summer}



\begin{quote}
Before the new one: say the Russian for winter, then the Russian for spring.
\end{quote}

\subsection*{You'll want to know: \ru{лето}}
\textbf{\ru{лето}} (\emph{léto}) — \textquotedblleft{}summer\textquotedblright{}.

\begin{cousinweb}
\textbf{\ru{лето}} $\leftarrow$ Proto-Slavic \emph{*lěto}, \textquotedblleft{}summer; year\textquotedblright{} — inherited. Its deeper origin is uncertain: Vasmer reports a possible link with Old Irish \emph{laithe}, \textquotedblleft{}day\textquotedblright{}, and Old Gutnish \emph{laþigs}, \textquotedblleft{}in spring\textquotedblright{}, and others with \emph{\ru{лить}}, \textquotedblleft{}to pour\textquotedblright{}. It also counts years: \emph{\ru{пять} \ru{лет}}, \textquotedblleft{}five years\textquotedblright{}.
\end{cousinweb}

\begin{grammarlens}[title={What you've built}]
One more word for this chapter.
\end{grammarlens}

\subsection*{Your turn}
\begin{itemize}
\raggedright
  \item \emph{Say it:} \emph{léto}
  \item \emph{Say it:} \emph{léto}, once more
  \item \emph{Say it:} say \emph{léto}
  \item \emph{Recall:} say the Russian for winter, then the Russian for spring, then say \emph{léto} again
\end{itemize}

\begin{tcolorbox}[breakable,colback=teal!4,colframe=teal!35!black,title={Before you move on}]
What does \ru{лето} mean? (\textquotedblleft{}Summer\textquotedblright{}.) Say it once more.
\end{tcolorbox}
\section[\texorpdfstring{\ru{осень} (ósen)}{осень (ósen)}]{\ru{осень} (ósen) — autumn}

\label{lesson:RU-C71-autumn}



\begin{quote}
Before the new one: say the Russian for spring, then the Russian for summer.
\end{quote}

\subsection*{You'll want to know: \ru{осень}}
\textbf{\ru{осень}} (\emph{ósen}) — \textquotedblleft{}autumn\textquotedblright{}.

\begin{cousinweb}
\textbf{\ru{осень}} $\leftarrow$ Proto-Slavic \emph{*esenĭ} — inherited; East Slavic turned \emph{je-} into \emph{\ru{о}-}. Vasmer relates it to Old Prussian \emph{assanis}, \textquotedblleft{}autumn\textquotedblright{}, Gothic \emph{asans}, \textquotedblleft{}harvest\textquotedblright{}, and Old High German \emph{aran}, \textquotedblleft{}harvest\textquotedblright{}: autumn as the reaping season.
\end{cousinweb}

\begin{grammarlens}[title={What you've built}]
One more word for this chapter.
\end{grammarlens}

\subsection*{Your turn}
\begin{itemize}
\raggedright
  \item \emph{Say it:} \emph{ósen}
  \item \emph{Say it:} \emph{ósen}, once more
  \item \emph{Say it:} say \emph{ósen}
  \item \emph{Recall:} say the Russian for spring, then the Russian for summer, then say \emph{ósen} again
\end{itemize}

\begin{tcolorbox}[breakable,colback=teal!4,colframe=teal!35!black,title={Before you move on}]
What does \ru{осень} mean? (\textquotedblleft{}Autumn\textquotedblright{}.) Say it once more.
\end{tcolorbox}
